\section{Separability-first design}\label{sec:design}
\paragraph{Margin and gap.} We use a margin $\delta=\DeltaNats$ nats per forecast, fixed before any result in the first
study and kept since. For a pair of references $P$ and $Q$ on a set of evaluation systems, the gap is the expected
log-score difference $\mathbb E[\log p(y\mid\text{context})-\log q(y\mid\text{context})]$ over contexts and targets from
those systems. When $P$ is the predictive of the process that generated the systems (population Bayes on held-out systems,
library Bayes on in-library systems), the gap equals an expected Kullback--Leibler divergence and is non-negative. On
designed sets (probe grids, out-of-support systems) neither reference is the truth, the gap can have either sign, and it
measures a behavioural difference rather than a divergence; we therefore report designed sets separately and do not use
them for the headline statements below.

\paragraph{Rules.} The studies used different pre-registered decision rules with the same margin: Stage~0 classified each
regime by an interval rule against a flat-prior frontier; \sHalf used intervals of $\pm$\SfiveIntervalSE standard errors with two
replicates that had to agree; \sSix used a studentised max-$|t|$ bootstrap over all contrasts. The reference-only maps
(1A, F1, F3) mark a cell \textsc{identifiable} if estimate $-\,3\,$SE $\ge\delta$, \textsc{unidentifiable} if
estimate $+\,3\,$SE $<\delta$, and \textsc{borderline} otherwise; the F1 map applies this to the absolute gap (two-sided),
which only matters on designed sets. For libraries each map draws several independent libraries per size $K$ (more in the
1A and F1 maps than in F3), applies the cell rule per draw, and marks a cell \textsc{identifiable} or \textsc{unidentifiable}
only if at least nine in ten draws agree, otherwise \textsc{partial}. Because a trained model sees one library, a
\textsc{partial} cell means that some libraries are separable from the continuous prior and others are not; we report
per-draw agreement rather than hide it. The maps use common random numbers, so their cells are dependent and patterns across
cells are descriptive; the rule is applied per cell without a global correction, and the maps guide design rather than test
hypotheses. Each map's specification was committed before its seed was first used.

\paragraph{What ``identifiable'' does and does not mean.} Here it is a property of two reference predictors at a
precision: their expected gap exceeds $\delta$ by more than three standard errors. It is necessary in our designs for a
per-cell behavioural classification at margin $\delta$, but not sufficient: a trained model is a random object whose seed
variance enters any classification. For F1, the bootstrap-calibrated operating characteristics of our pre-registered
classification rule imply that, with \MinModels models per arm, a population-like versus context-only classification has
at least eighty per cent power only if the between-seed standard deviation of a model's regret is at most
\LookupQoneThirtyTwo nats (\LookupQoneSixtyFour at twice as many models). Evidence pooled over many cells, for example
by a sequence-level test, could also separate references whose per-cell gaps are below $\delta$; our maps do not use it.
